\chapter{一般拓扑}\label{ch:b3:topology}

本科二年级分册讨论了度量空间中的分析：距离、
球、数列。但基本概念——连续性、
紧性、连通性——从不涉及距离的具体数值，
只关心它所生成的\emph{\omterm{def:b3:topology:topology}{开集}}族。
本章以这一族集合为基本对象。这样做的收益
不在于一般性本身：商构造（
圆周 $\R/\Z$、射影空间）、乘积，以及
泛函分析中的弱型\omterm{def:b3:topology:topology}{拓扑}等，本就不是
以度量为先的对象。先重新建立连续性，再
用开覆盖讨论紧性（并证明它与本科二年级分册中
度量空间里的序列定义等价），以及连通性——
最后以一条定理作结：连续双射不能将
$\R$ 与 $\R^2$ 等同起来：\omterm{def:b3:topology:topology}{拓扑}能够区分维数。

\section{拓扑、开集与连续性}

\begin{definition}\label{def:b3:topology:topology}
\emph{拓扑}\index{拓扑}是集合 $X$ 上的一个
子集族 $\mathcal T$（$X$ 的子集）——其元素称为\emph{开
集}\index{开集}——满足：$\varnothing, X \in
\mathcal T$；任意开集之并仍为开集；任意\emph{有限}
个开集之交仍为开集。$(X, \mathcal T)$ 称为
\emph{拓扑空间}。开集的补集称为
\emph{闭集}。\emph{邻域}\index{邻域}：对 $x$ 而言，
是包含某个含 $x$ 的开集的集合。
\end{definition}

\begin{example}\label{ex:b3:topology:examples}
(a) 度量空间，“开”的含义同本科二年级分册（开
球之并）：\emph{度量\omterm{def:b3:topology:topology}{拓扑}}；不同的度量可以给出同一\omterm{def:b3:topology:topology}{拓扑}
（等价度量）。若\omterm{def:b3:topology:topology}{拓扑}由某个度量给出，
则称该空间为\emph{可度量化}的。
(b) \emph{离散}\omterm{def:b3:topology:topology}{拓扑}（即全部子集）与
\emph{平凡}\omterm{def:b3:topology:topology}{拓扑} $\{\varnothing, X\}$。
(c) 无限集上的\emph{余有限}\omterm{def:b3:topology:topology}{拓扑}：\omterm{def:b3:topology:topology}{开集} $=$
空集或余有限集。该\omterm{def:b3:topology:topology}{拓扑}不可度量化，
（见\cref{exo:b3:topology:3}）。
(d) 在 $\R$ 上取通常\omterm{def:b3:topology:topology}{拓扑}；在 $\bar\R = \R \cup
\{\pm\infty\}$ 上由射线生成的序\omterm{def:b3:topology:topology}{拓扑}——使得
“$x_n \to +\infty$”成为普通收敛的一个实例。
\end{example}

\begin{definition}\label{def:b3:topology:interior}
对 $A \subseteq X$，\emph{内部} $\mathring A$ 为
包含于 $A$ 中的最大\omterm{def:b3:topology:topology}{开集}（即所有这样的\omterm{def:b3:topology:topology}{开集}之并）；
\emph{闭包}\index{闭包} $\bar A$ 为最小的闭集，
它包含 $A$；\emph{边界} $\partial A = \bar A
\setminus \mathring A$。$A$ 称为\emph{稠密集}\index{稠密集}，若
$\bar A = X$。易知，$x \in \bar A$ 当且仅当
$x$ 的每个\omterm{def:b3:topology:topology}{邻域}都与 $A$ 相交（若某一\omterm{def:b3:topology:topology}{邻域}与 $A$ 不交，则其
内部的补集是包含 $A$ 的更小闭集；反之亦然）。
\end{definition}

\begin{definition}\label{def:b3:topology:basis}
\emph{拓扑基}\index{拓扑基}是 $\mathcal T$ 的
一族子集 $\mathcal B \subseteq \mathcal T$，满足：每个\omterm{def:b3:topology:topology}{开集}
都是 $\mathcal B$ 中成员之并（例如度量空间中的开球
，或 $\R$ 中的开区间）。$\mathcal B$
为 $X$ 的子集族，它为\emph{某个}\omterm{def:b3:topology:topology}{拓扑}的\omterm{def:b3:topology:topology}{拓扑}基，当且仅当它覆盖
$X$，并且对于 $B_1, B_2 \in \mathcal B$ 及 $x \in B_1 \cap B_2$，
存在 $B_3 \in \mathcal B$ 使得 $x \in B_3 \subseteq B_1\cap B_2$，
则“成员之并”构成一个\omterm{def:b3:topology:topology}{拓扑}，即
由 $\mathcal B$ \emph{生成}的\omterm{def:b3:topology:topology}{拓扑}。
\end{definition}

\begin{definition}\label{def:b3:topology:continuity}
映射 $f \colon X \to Y$ 称为\emph{连续映射}\index{连续映射}，若
$f^{-1}(V)$ 为\omterm{def:b3:topology:topology}{开集}，对每个\omterm{def:b3:topology:topology}{开集} $V \subseteq Y$——
等价地，闭集的原像为闭集；等价地，
对每个 $x$，若 $V$ 为 $f(x)$ 的\omterm{def:b3:topology:topology}{邻域}，则 $f^{-1}(V)$
为 $x$ 的\omterm{def:b3:topology:topology}{邻域}（即在每一点 $x$ 处\emph{连续}）。
只需在 $Y$ 的\omterm{def:b3:topology:basis}{拓扑基}上验证即可。连续映射的复合
仍连续。\emph{同胚}\index{同胚}是
一个逆映射也连续的连续双射；\omterm{def:b3:topology:topology}{拓扑}学研究
在同胚下保持不变的性质。
\end{definition}

\begin{proposition}[连续性与闭包；粘合]
\label{prop:b3:topology:gluing}
(a) $f$ 连续，当且仅当 $f(\bar A) \subseteq \overline{f(A)}$
对所有 $A \subseteq X$ 成立。
(b) 若 $X = F_1 \cup F_2$，其中 $F_1, F_2$ 为闭集，且 $f\colon X
\to Y$ 在每个 $F_i$ 上的限制都连续，则 $f$
连续。
\end{proposition}

\begin{proof}
(a) 若 $f$ 连续，则 $f^{-1}(\overline{f(A)})$ 为闭集，且
包含 $A$，从而包含 $\bar A$。反之，
对 $A = f^{-1}(F)$（$F$ 为闭集）应用上述判别准则：$f(\bar A) \subseteq
\overline{f(f^{-1}(F))} \subseteq \bar F = F$，于是 $\bar A
\subseteq f^{-1}(F) = A$，即闭集的原像为闭集。
(b) 对闭集 $F \subseteq Y$，有 $f^{-1}(F) =
(f\restriction_{F_1})^{-1}(F) \cup (f\restriction_{F_2})^{-1}
(F)$，是分别在 $F_1$、$F_2$ 中为闭的两个集合之并，从而
在 $X$ 中为闭（$F_i$ 为闭集，闭集的闭子集仍为闭集）。
\end{proof}

\begin{definition}\label{def:b3:topology:hausdorff}
$X$ 称为\emph{豪斯多夫空间}\index{豪斯多夫空间}（亦称
\emph{分离}空间），若任意两个不同点都有不交的
\omterm{def:b3:topology:topology}{邻域}。度量空间都是豪斯多夫的（取半径为
$d(x,y)/2$ 的球）。序列 $(x_n)$ \emph{收敛}于 $x$，若
$x$ 的每个\omterm{def:b3:topology:topology}{邻域}都含有除有限个之外的所有 $x_n$；在
豪斯多夫空间中，极限是唯一的（若有两个极限，
则二者各有不交的\omterm{def:b3:topology:topology}{邻域}，且每个\omterm{def:b3:topology:topology}{邻域}都含某个尾部）。在
豪斯多夫空间中，点——从而有限集——都是闭集。
\end{definition}

\begin{remark}\label{rem:b3:topology:sequences}
在度量空间中，序列能刻画一切：$x \in \bar A$
当且仅当 $A$ 中存在某个序列收敛于 $x$（取 $x_n \in A \cap
B(x, 1/n)$），且 $f$ 连续当且仅当它
序列连续。在一般空间中两个等价性均不成立；
正确的序列替代物（滤子、网）属于更高阶的
课程。在度量空间中，我们将陈述基于序列的结论，
在一般空间中则陈述基于覆盖的结论——并在适当处
证明二者等价。
\end{remark}

\section{子空间、乘积、商}

\begin{definition}\label{def:b3:topology:constructions}
由旧空间构造新空间的三种方式：\index{乘积
拓扑}\index{商拓扑}
\begin{enumerate}
  \item \emph{子空间}：在 $A \subseteq X$ 上，\omterm{def:b3:topology:topology}{开集}为
        $U \cap A$（$U$ 为 $X$ 中的\omterm{def:b3:topology:topology}{开集}）——
        使包含映射连续的最粗\omterm{def:b3:topology:topology}{拓扑}。
  \item \emph{乘积}：在 $X \times Y$（及有限乘积）上，
        \omterm{def:b3:topology:topology}{拓扑}以\emph{开矩形} $U \times
        V$ 为基；对无限乘积 $\prod_i X_i$，基
        由形如 $\prod U_i$ 的盒子组成，其中 $U_i = X_i$ 对
        \emph{除有限个外的所有} $i$ ——使每个投影连续的最粗
        \omterm{def:b3:topology:topology}{拓扑}。
  \item \emph{商}：若 $\sim$ 为 $X$ 上的等价关系，且
        $\pi\colon X \to X/{\sim}$ 为投影，规定 $V
        \subseteq X/{\sim}$ 为\omterm{def:b3:topology:topology}{开集}，当且仅当 $\pi^{-1}(V)$ 为\omterm{def:b3:topology:topology}{开集}——
        使 $\pi$ 连续的最细\omterm{def:b3:topology:topology}{拓扑}。
\end{enumerate}
\end{definition}

\begin{proposition}[泛性质]
\label{prop:b3:topology:universal}
(a) $f \colon Z \to X \times Y$ 连续当且仅当它的
两个分量 $\pi_X \circ f$、$\pi_Y \circ f$ 均连续
（对任意乘积同理）。
(b) $g \colon X/{\sim} \to Z$ 连续当且仅当 $g \circ \pi
\colon X \to Z$ 连续。
\end{proposition}

\begin{proof}
(a) 必要性：由复合即得。充分性：只需检验
基中开矩形的原像：$f^{-1}(U \times V) = (\pi_Xf)^{-1}(U)
\cap (\pi_Yf)^{-1}(V)$，为\omterm{def:b3:topology:topology}{开集}（无限情形下仅有限个因子 $\neq X_i$
）。(b) 必要性：由复合即得。充分性：
设 $W \subseteq Z$ 为\omterm{def:b3:topology:topology}{开集}，则 $\pi^{-1}(g^{-1}(W)) = (g\pi)^{-1}(W)$
是\omterm{def:b3:topology:topology}{开集}，由\omterm{def:b3:topology:constructions}{商拓扑}的定义，这意味着
$g^{-1}(W)$ 是\omterm{def:b3:topology:topology}{开集}。
\end{proof}

\begin{example}\label{ex:b3:topology:circle}
商空间 $\R/\Z$（等同 $x$ 与 $x + n$）\omterm{def:b3:topology:continuity}{同胚}于
圆周 $S^1 = \{z \in \C : \abs z = 1\}$：映射 $x
\mapsto \eu^{2\iu\pi x}$ 诱导出一个连续双射
$\R/\Z \to S^1$（由\cref{prop:b3:topology:universal}(b)）；其
逆映射由紧性论证知其连续，该论证见下文
\cref{cor:b3:topology:compacthomeo}。
（\cref{exo:b3:topology:5} 给出全部细节，包括
$\R/\Z$ 为何是豪斯多夫的且紧的）。同理，$[0,1]$ 的两端
点粘合后即为 $S^1$，正方形对边粘合后
得环面；粘合归根结底是一个定理，而非图画。
\end{example}

\section{紧性}

\begin{definition}\label{def:b3:topology:compact}
\emph{开覆盖}是 $X$ 的一族\omterm{def:b3:topology:topology}{开集} $(U_i)_{i\in I}$，
满足 $\bigcup U_i = X$。称 $X$ 为
\emph{紧}\index{紧空间}，若它是豪斯多夫的，且每个
开覆盖都存在\emph{有限子覆盖}。等价地（取补集）：
每个具有\emph{有限
交性质}的闭集族，即任意有限子族的交非空，
其总交亦非空。
\end{definition}

\begin{theorem}[基本性质]\label{thm:b3:topology:compactprops}
设 $X$ 为\omterm{def:b3:topology:compact}{紧空间}。
\begin{enumerate}
  \item $X$ 的闭子集是紧的；\omterm{def:b3:topology:hausdorff}{豪斯多夫空间}的紧子集
        是闭集。
  \item \omterm{def:b3:topology:compact}{紧空间}在\omterm{def:b3:topology:hausdorff}{豪斯多夫空间}中的连续像
        是紧的。特别地，\omterm{def:b3:topology:continuity}{连续映射} $f\colon X
        \to \R$ 有界，且可取到其上下确界。
  \item $X$ 中非空闭子集构成的递减序列
        的交非空。
\end{enumerate}
\end{theorem}

\begin{proof}
(1) 设 $F \subseteq X$ 为闭集，$(U_i)$ 为 $F$ 的开覆盖
（由 $X$ 的\omterm{def:b3:topology:topology}{开集}组成）：添加 $X \setminus F$ 即得 $X$ 的开覆盖，
有限子覆盖去掉 $X \setminus F$ 后仍覆盖 $F$。
子空间的豪斯多夫性显然。反之，设 $K \subseteq Y$
为紧集，$Y$ 为豪斯多夫的，且 $y \notin K$：对每个 $x \in K$
取不相交的\omterm{def:b3:topology:topology}{开集} $U_x \ni x$、$V_x \ni y$；有限多个 $U_x$
覆盖 $K$，相应的 $V_x$ 之交为 $y$ 的
\omterm{def:b3:topology:topology}{邻域}，与 $K$ 的覆盖不交，从而与 $K$ 不交：
即 $K$ 的补集是\omterm{def:b3:topology:topology}{开集}。

(2) 若 $(V_j)$ 覆盖 $f(X)$，则 $(f^{-1}(V_j))$ 覆盖 $X$；
像空间中的有限子覆盖来自原空间中的有限子覆盖。
像作为子空间是豪斯多夫的。对实值 $f$：
$f(X)$ 在 $\R$ 中紧，从而既闭又有界（用
$(-n, n)$ 覆盖证有界；由 (1) 得闭），且闭有界集
必含其上确界。

(3) 若 $\bigcap F_n = \varnothing$，则\omterm{def:b3:topology:topology}{开集} $X \setminus F_n$
覆盖 $X$；有限个即可，故存在 $F_{n_0} =
\varnothing$（因序列递减）——矛盾。
\end{proof}

\begin{corollary}\label{cor:b3:topology:compacthomeo}
从\omterm{def:b3:topology:compact}{紧空间}到\omterm{def:b3:topology:hausdorff}{豪斯多夫空间}的连续\emph{双射}
是\omterm{def:b3:topology:continuity}{同胚}。
\end{corollary}

\begin{proof}
逆映射连续，当且仅当闭集的像都是闭集；
取闭集 $F$，它是紧的
（\cref{thm:b3:topology:compactprops}(1)），其像是紧的
（2），从而在豪斯多夫的目标空间中是闭的（1）。
\end{proof}

\begin{theorem}[有限积]\label{thm:b3:topology:tychonoff}
有限个\omterm{def:b3:topology:compact}{紧空间}的积是紧的。\index{吉洪诺夫
定理（有限情形）}
\end{theorem}

\begin{proof}
只需讨论 $X \times Y$。豪斯多夫性可遗传
（在某一个坐标上分离即可）。设 $(W_i)$ 为 $X
\times Y$ 的开覆盖；不妨设各 $W_i$ 为\omterm{def:b3:topology:basis}{拓扑基}中的方块 $U_i \times
V_i$（加细：每点都位于含于某个 $W_i$ 的方块中；方块的
有限子覆盖可得 $W_i$ 的一个有限子覆盖）。固定 $x \in X$：
\emph{切片} $\{x\}\times Y \cong Y$ 是紧的，故有限
个方块 $U_1\times V_1, \dots, U_k\times V_k$ 覆盖它，且
$x \in U_j$ 对所有 $j$ 成立；于是 $U^x = \bigcap_{j \leq k} U_j$ 为
$x$ 的开\omterm{def:b3:topology:topology}{邻域}，且 $U^x \times Y$ 被
有限个方块覆盖（\emph{管引理}\index{管引理}：
对 $(x', y) \in U^x \times Y$，$y \in V_j$
对某个 $j$ 成立，因为这些方块覆盖了 $\{x\}\times Y$ 在水平 $y$ 处的部分，
又 $x' \in U^x \subseteq U_j$，故 $(x', y) \in U_j
\times V_j$）。现有限个 $U^x$
覆盖\omterm{def:b3:topology:compact}{紧空间} $X$；相应的各有限方块族
合起来覆盖 $X \times Y$。
\end{proof}

\begin{theorem}[度量空间中的紧性]
\label{thm:b3:topology:metriccompact}
设 $(X, d)$ 为度量空间，下列三条件
等价：\index{序列紧性}
\begin{enumerate}
  \item $X$ 是紧的（博雷尔–勒贝格）；
  \item $X$ 中每个序列都有收敛子序列；
        （\omterm{thm:b3:topology:metriccompact}{序列紧性}——本科二年级的定义）；
  \item $X$ 是完备的，且\emph{全有界}：对每个
        $\varepsilon > 0$，存在有限个半径为
        $\varepsilon$ 的球覆盖 $X$。
\end{enumerate}
\end{theorem}

\begin{proof}
(1)$\Rightarrow$(2)：设 $(x_n)$ 没有收敛
子序列。于是每个 $x \in X$ 都有开球 $B_x$
使 $x_n$ 仅对有限多个 $n$ 落在其中（否则
存在子序列收敛到 $x$，取半径 $1/k$ 的球）。有限个
$B_x$ 覆盖 $X$，故下标 $n$ 仅有有限个，
矛盾。

(2)$\Rightarrow$(3)：完备性：带收敛子序列的柯西序列
必收敛（本科二年级内容）。全有界性：若
某个 $\varepsilon$ 没有有限覆盖，则归纳选取
$x_{n+1}$ 位于 $\bigcup_{k\leq n} B(x_k, \varepsilon)$ 之外。此时
该序列中 $d(x_m, x_n) \geq \varepsilon$ 对 $m \neq n$ 成立，
故既无柯西子序列，也无收敛子序列。

(3)$\Rightarrow$(1)：先证 (3) 蕴含 (2)：给定 $(x_n)$，
$X$ 可由有限个半径为 $1$ 的球覆盖，其中之一 $B_1$，
含有一个子序列；$X$ 可由半径为 $1/2$ 的球覆盖，其中之一
含有更进一步的子序列；迭代并对角化，其
对角子序列是柯西序列（第 $k$ 步以后的两项落在
半径为 $2^{-k}$ 的同一个球中，在通常的 $2\cdot$ 意义下）。
于是收敛。现设 $(U_i)$ 为开覆盖，并假设没有
有限子覆盖。\emph{勒贝格数论证}：对每个 $n$，
存在球 $B(y_n, 2^{-n})$ 不能被有限个
$U_i$ 覆盖——事实上，$X$ 可由有限个半径为
$2^{-n}$ 的球覆盖；若每个球都可被有限个覆盖，则 $X$ 亦然。由
(2)，存在子序列 $y_{n_k} \to y$；取 $i$ 使 $y \in U_i$
且 $r > 0$ 使 $B(y, r) \subseteq U_i$。当 $k$ 充分大时，
$B(y_{n_k}, 2^{-n_k}) \subseteq B(y, r) \subseteq U_i$，被
\emph{一个} $U_i$ 覆盖——矛盾。
\end{proof}

\begin{corollary}[海涅–博雷尔；海涅]
\label{cor:b3:topology:heineborel}
(a) $\R^n$ 的子集是紧的当且仅当它既闭又有界。
(b) 紧度量空间到度量空间的\omterm{def:b3:topology:continuity}{连续映射}
是一致连续的。\index{海涅–博雷尔
定理}\index{海涅定理}
\end{corollary}

\begin{proof}
(a) 闭且有界 $\Rightarrow$ 含于某个立方体
$[-M,M]^n$ 是紧的：$[-M, M]$ 紧（序列意义下，由
波尔查诺–魏尔斯特拉斯——或直接用覆盖的二分法），
由\cref{thm:b3:topology:tychonoff}可得乘积的紧性；然后
应用 \cref{thm:b3:topology:compactprops}(1) 即得。反之，
紧子集是闭的
（\cref{thm:b3:topology:compactprops}(1)），且有界（用
同心球覆盖）。

(b) 设 $f \colon X \to Y$，$\varepsilon > 0$。
以下各球 $B\bigl(x, \delta_x\bigr)$ 满足 $f\bigl(B(x,
2\delta_x)\bigr)\subseteq B\bigl(f(x), \varepsilon/2\bigr)$，
这些球覆盖 $X$；取有限子覆盖 $B(x_j, \delta_{x_j})$，并
令 $\delta = \min_j \delta_{x_j}$。若 $d(x, x') < \delta$：$x
\in B(x_j, \delta_{x_j})$ 对某个 $j$ 成立，于是二者 $x, x'
\in B(x_j, 2\delta_{x_j})$，故 $d(f(x), f(x')) < \varepsilon$。
\end{proof}

\begin{definition}\label{def:b3:topology:locallycompact}
称 $X$ 为\emph{局部紧空间}\index{局部紧空间}，若它
是豪斯多夫的，且每点都有紧\omterm{def:b3:topology:topology}{邻域}（$\R^n$；
$\R^n$ 的开子集；离散空间——但 $\Q$ 不是，见
\cref{exo:b3:topology:9}）。每个局部\omterm{def:b3:topology:compact}{紧空间}都可嵌入一个\omterm{def:b3:topology:compact}{紧空间}，
即\emph{单点
紧化}\index{单点紧化} $\hat X = X
\cup \{\infty\}$，其\omterm{def:b3:topology:topology}{开集}是 $X$ 的\omterm{def:b3:topology:topology}{开集}，连同
在 $\hat X$ 中，$X$ 的紧子集的补集。
直接验证公理；$\hat X$ 是紧的（覆盖中必有一个成员
含 $\infty$，其补集是紧的，被
有限个其他成员覆盖），且是豪斯多夫的（将 $x$ 与 $\infty$
用 $x$ 的紧\omterm{def:b3:topology:topology}{邻域}及其补集分开）。例：
$\hat\R \cong S^1$，且 $\widehat{\R^n} \cong S^n$ 由
立体投影得到（见\cref{exo:b3:topology:11}）。
\end{definition}

\section{连通性}

\begin{definition}\label{def:b3:topology:connected}
称 $X$ 为\emph{连通空间}\index{连通空间}，若它不是
两个不相交的非空\omterm{def:b3:topology:topology}{开集}之并——等价地，其
既开又闭的子集只有 $\varnothing$ 与 $X$；
等价地，每个\omterm{def:b3:topology:continuity}{连续映射} $X \to \{0, 1\}$（离散）
都是常值的。子集连通，是指它作为子空间连通。
\end{definition}

\begin{theorem}\label{thm:b3:topology:connectedbasics}
\begin{enumerate}
  \item $\R$ 的连通子集恰为区间。
  \item 连通集的连续像仍连通（由此可得
        介值定理：\omterm{def:b3:topology:connected}{连通空间}上的连续实映射
        的像是一个区间）。
  \item 若 $(A_i)$ 均连通且有公共点，则 $\bigcup
        A_i$ 连通。若 $A$ 连通且 $A \subseteq
        B \subseteq \bar A$，则 $B$ 连通。
  \item 有限个\omterm{def:b3:topology:connected}{连通空间}的积是连通的。
\end{enumerate}
\end{theorem}

\begin{proof}
全文使用 $\{0,1\}$ 判别法：\omterm{def:b3:topology:continuity}{连续映射} $f \colon X \to
\{0,1\}$ 必为常值。

（1）若 $A$ 不是区间，则存在 $c$ 介于 $a, b \in A$ 之间而不属于它：
$A = (A \cap (-\infty, c)) \sqcup (A \cap (c, +\infty))$
于是该集合不连通。反之，设 $I$ 为区间，$f \colon I
\to \{0,1\}$ 连续，且 $f(a) = 0$，$f(b) = 1$，$a < b$。
令 $c = \sup\{x \in [a,b] : f(x) = 0\}$；由 $c$ 处的连续性
可得 $f(c) = 0$（取值 $0$ 的极限：$c$ 的每个\omterm{def:b3:topology:topology}{邻域}
都与 $\{f = 0\}$ 相交；$\{f=0\}$ 是闭的），且 $c < b$
且 $f \equiv 1$ 于 $(c, b]$ 上，故 $f(c) = 1$（由
$\{f = 1\} \ni c$ 上同样的闭性论证），矛盾。

（2）\omterm{def:b3:topology:continuity}{连续映射} $g \colon f(X) \to \{0,1\}$ 给出 $g \circ f$
为常值，故 $g$ 在像集上为常值。

（3）\omterm{def:b3:topology:continuity}{连续映射} $f\colon \bigcup A_i \to \{0,1\}$ 在每个 $A_i$ 上为常值，
且在公共点处取值相同。关于
\omterm{def:b3:topology:interior}{闭包}：$f \colon B \to \{0,1\}$ 为常值 $= c$（于 $A$ 上）；任取
$x \in B \subseteq \bar A$，属于 $A$ 的\omterm{def:b3:topology:interior}{闭包}，且 $f^{-1}
(f(x))$ 是 $x$ 的\omterm{def:b3:topology:topology}{邻域}（\omterm{def:b3:topology:topology}{开集}的原像），它必
与 $A$ 相交，故 $f(x) = c$。

（4）对 $X \times Y$ 及 $f\colon X\times Y \to \{0,1\}$，任取
两点 $(x,y), (x',y')$ 可由“折线”
$\{x\}\times Y \cup X \times \{y'\}$ 相连，它是两个连通集
之并（分别与 $Y$、$X$ \omterm{def:b3:topology:continuity}{同胚}），两者交于 $(x, y')$。由（3）及
（2），$f$ 在两点处取值相同。
\end{proof}

\begin{definition}\label{def:b3:topology:pathconnected}
称 $X$ 为\emph{道路连通空间}\index{道路连通空间}，若任意
两点可由一条\emph{道路}（\omterm{def:b3:topology:continuity}{连续映射} $\gamma
\colon [0,1] \to X$）相连。道路连通蕴含连通：
对\omterm{def:b3:topology:continuity}{连续映射} $f \colon X \to \{0,1\}$，$x, y$ 处的两个取值
都是常值映射 $f\circ\gamma$ 的取值，故相等
（\cref{thm:b3:topology:connectedbasics}(1)--(2)）。赋范空间中的凸子集
道路连通（用线段连接）；$\R^n
\setminus \{0\}$（$n \geq 2$）亦然（绕原点走），以及
$S^n$（$n \geq 1$）亦然（将 $\R^{n+1}\setminus
\{0\}$ 中的道路投影而得）。
\end{definition}

\begin{example}[拓扑学家的正弦曲线]
\label{ex:b3:topology:sinecurve}
设 $\Gamma = \{(x, \sin\frac1x) : 0 < x \leq 1\}$，$S =
\bar\Gamma = \Gamma \cup (\{0\}\times[-1,1])$（每个点 $(0,
y)$，$\abs y \leq 1$，都是 $\Gamma$ 中点的极限：解
$\sin\frac1x = y$（在 $0$ 附近）。于是 $S$
\emph{连通}——连通集 $\Gamma$ 的\omterm{def:b3:topology:interior}{闭包}，
是 $(0, 1]$ 的连续像
（\cref{thm:b3:topology:connectedbasics}(3)）——但\emph{不是
\omterm{def:b3:topology:pathconnected}{道路}连通}：从 $(1, \sin 1)$ 到 $(0,0)$ 的\omterm{def:b3:topology:pathconnected}{道路}
必须让横坐标 $\to 0$，而纵坐标振荡
于 $\pm1$ 之间；此点可由\cref{exo:b3:topology:9}严格证明。
连通性与\omterm{def:b3:topology:pathconnected}{道路}连通性确实
不同。\index{拓扑学家的正弦曲线}
\end{example}

\begin{omfigure}
\begin{tikzpicture}
\begin{axis}[omaxis, width=10.5cm, height=4.6cm,
    xmin=-0.02, xmax=0.55, ymin=-1.25, ymax=1.25,
    xtick={0, 0.25, 0.5}, ytick={-1, 1}]
  \addplot[omThm, thick, domain=0.006:0.55, samples=1200]
    {sin(deg(1/x))};
  \addplot[omDef, very thick] coordinates {(0,-1) (0,1)};
\end{axis}
\end{tikzpicture}

{\small \omterm{ex:b3:topology:sinecurve}{拓扑学家的正弦曲线}：函数
$\sin\frac1x$ 的图像（红色）在整个线段上聚集
$\{0\}\times[-1,1]$（蓝色）。并集连通，但不是
\omterm{def:b3:topology:pathconnected}{道路}连通的：没有连续\omterm{def:b3:topology:pathconnected}{道路}能在有限的参数时间内
穿越无穷多次振荡。}
\end{omfigure}

\begin{definition}\label{def:b3:topology:components}
\emph{连通分支}\index{连通分支}是指：对 $x
\in X$，它是所有包含 $x$ 的连通子集之并——
即其中最大者（\cref{thm:b3:topology:connectedbasics}(3)）。
分支构成 $X$ 的一个划分，且为闭集（连通集的\omterm{def:b3:topology:interior}{闭包}
仍连通）。称 $X$ 为\emph{完全不连通的}，若
其所有分支皆为单点（$\Q$；周末
问题中的康托尔集）。
\end{definition}

\begin{theorem}\label{thm:b3:topology:rnotr2}
$\R$ 与任何 $\R^n$（$n \geq 2$）都不\omterm{def:b3:topology:continuity}{同胚}。
\end{theorem}

\begin{proof}
设 $h \colon \R \to \R^n$ 为\omterm{def:b3:topology:continuity}{同胚}。去掉一点后，
$\R \setminus \{0\}$ 与 $\R^n \setminus
\{h(0)\}$ \omterm{def:b3:topology:continuity}{同胚}。但 $\R\setminus\{0\}$ 不连通，而
$\R^n\setminus\{\text{点}\}$ 当 $n \geq
2$ 时\omterm{def:b3:topology:pathconnected}{道路}连通（见\cref{def:b3:topology:pathconnected}；将该点平移到
$0$）。连通性是\omterm{def:b3:topology:continuity}{同胚}不变量
（见\cref{thm:b3:topology:connectedbasics}(2)）——矛盾。
（$\R^2 \not\cong \R^3$ 需要更精细的不变量——
代数\omterm{def:b3:topology:topology}{拓扑}；周末问题提醒我们注意其中的陷阱：
\emph{连续满射} $\R \to \R^2$ 确实存在。）
\end{proof}

\begin{method}\label{met:b3:topology:connectargument}
证明一个集合连通：可将其表示为连续像，
相交连通集之并、\omterm{def:b3:topology:interior}{闭包}，或乘积；
（见\cref{thm:b3:topology:connectedbasics}）；对赋范空间的子集，
证明其\omterm{def:b3:topology:pathconnected}{道路}连通时可用显式\omterm{def:b3:topology:pathconnected}{道路}（线段、
弧、折线）。要证明两个空间\emph{不}\omterm{def:b3:topology:continuity}{同胚}，
只需找到一个取值不同的\omterm{def:b3:topology:topology}{拓扑}不变量——紧性、
连通性、\omterm{def:b3:topology:components}{连通分支}数，或适当选定的有限集被删去后的各\omterm{def:b3:topology:components}{连通分支}
（删去有限集的技巧，见
\cref{thm:b3:topology:rnotr2}：它还证明了 $[0,1) \not\cong
(0,1)$ 与 $S^1 \not\cong [0,1]$）。
\end{method}

\section{练习}

\begin{exercise}[$\star$]\label{exo:b3:topology:1}
(a) 列出 $\{a, b\}$ 上的全部\omterm{def:b3:topology:topology}{拓扑}，在\omterm{def:b3:topology:continuity}{同胚}意义下对其分类，
并确定其中哪些是连通的、哪些是
豪斯多夫的。
(b) 在 $\R$（通常\omterm{def:b3:topology:topology}{拓扑}）中，求下列集合的\omterm{def:b3:topology:interior}{内部}、\omterm{def:b3:topology:interior}{闭包}与边界：
$\Q$、$[0,1) \cup \{2\}$ 与 $\{1/n : n \geq
1\}$。
\end{exercise}

\begin{exercise}[$\star$]\label{exo:b3:topology:2}
(a) 证明 $f\colon X \to Y$ 连续，当且仅当
$f^{-1}(B)$ 对 $B$ 取遍 $Y$ 某一固定\omterm{def:b3:topology:basis}{拓扑基}的元素时均为\omterm{def:b3:topology:topology}{开集}。
(b) 证明 $f = \mathbf 1_{[0,\infty)} \colon \R \to \R$ 不连续，
但它是\emph{右连续}的。验证
半开区间 $[a, b)$ 构成
源空间 $\R$ 上某一\omterm{def:b3:topology:topology}{拓扑}的基（该\omterm{def:b3:topology:topology}{拓扑}称为\emph{索根弗雷直线}），
该\omterm{def:b3:topology:topology}{拓扑}严格细于通常\omterm{def:b3:topology:topology}{拓扑}，并证明：函数 $\R \to
\R$（目标空间取通常\omterm{def:b3:topology:topology}{拓扑}）自索根弗雷直线出发连续，当且仅当
它在每一点处都右连续。
\end{exercise}

\begin{exercise}[$\star$]\label{exo:b3:topology:3}
在无限集上赋予余有限\omterm{def:b3:topology:topology}{拓扑}，证明：任意两个
非空\omterm{def:b3:topology:topology}{开集}必相交（因此该空间不是豪斯多夫的，从而
不可度量化）；每个单射序列都收敛于
\emph{每一点}。极限唯一性的证明
在何处用到了豪斯多夫性？
\end{exercise}

\begin{exercise}[$\star\star$]\label{exo:b3:topology:4}
(a) 证明乘积的投影 $\pi_X, \pi_Y$ 连续且
\emph{开}（即\omterm{def:b3:topology:topology}{开集}的像是\omterm{def:b3:topology:topology}{开集}），但一般不是
闭映射（以 $\{xy = 1\}$ 在 $\R^2$ 中为例）。
(b) 证明可数个度量空间的乘积 $\prod_n X_n$ 中的序列
收敛当且仅当每个分量都收敛，并证明
$d(x, y) = \sum_n 2^{-n}\min(1, d_n(x_n, y_n))$ 度量化了
\omterm{def:b3:topology:constructions}{乘积拓扑}。
(c) 在 $\R^\N$ 上，证明“箱\omterm{def:b3:topology:topology}{拓扑}”（任意\omterm{def:b3:topology:topology}{开集}之积为\omterm{def:b3:topology:topology}{开集}）
严格细于\omterm{def:b3:topology:constructions}{乘积拓扑}：序列
$x^{(k)} = (1/k, 1/k, \dots)$ 按\omterm{def:b3:topology:constructions}{乘积拓扑}收敛于 $0$，
但按箱\omterm{def:b3:topology:topology}{拓扑}不收敛。
\end{exercise}

\begin{exercise}[$\star\star$]\label{exo:b3:topology:5}
圆周的三种看法。证明以下空间两两\omterm{def:b3:topology:continuity}{同胚}，
并给出显式映射：(i) $S^1 \subseteq \R^2$；
(ii) $\R/\Z$（\omterm{def:b3:topology:constructions}{商拓扑}）；(iii) $[0,1]/(0 \sim 1)$。
\emph{（对 (ii)：证明 $\R/\Z$ 是豪斯多夫的——取两个等价类，
使其代表元之间的距离 $\leq \frac12$——并证明
$\pi([0,1])$ 即为全空间，因此商空间是紧的；再用
\cref{cor:b3:topology:compacthomeo}。）}
\end{exercise}

\begin{exercise}[$\star\star$]\label{exo:b3:topology:6}
设 $X$ 为度量空间。
(a) 证明有限个紧子集之并是紧的，
且任意多个紧子集之交仍是紧的。
(b) 若 $K$ 紧、$F$ 闭，且 $K \cap F = \varnothing$，
证明 $d(K, F) = \inf\{d(x,y) : x \in K, y \in F\} > 0$；
并以两个不相交的闭集给出反例。
(c) 证明 $X$ 紧当且仅当\emph{每个}连续函数 $f
\colon X \to \R$ 有界。\emph{（若某序列没有收敛子列，则
可构造一个无界连续函数，
使其支撑在该序列各项附近；或利用形如 $x \mapsto
d(x, \cdot)$ 型函数——一条简洁的路线：若 $(x_n)$
没有聚点，则集合 $\{x_n\}$ 是闭的离散集，且
$f(x_n) = n$ 可借助不依赖蒂茨定理的方法连续延拓：
$f(x) = \sum_n n\,\max\bigl(0, 1 - \frac{d(x, x_n)}{r_n}\bigr)$
只需 $r_n$ 充分小。）}
\end{exercise}

\begin{exercise}[$\star\star$]\label{exo:b3:topology:7}
（勒贝格数）设 $(U_i)$ 是紧度量空间
$X$ 的一个开覆盖。证明存在 $\delta > 0$，使得每个
直径 $< \delta$ 的子集都含于某一个 $U_i$ 之中。
\emph{（否则可取 $A_n$，使其直径 $< 1/n$ 且不含于任何 $U_i$，
并取所选 $x_n \in A_n$ 的一个聚点。）} 由此再次推出
海涅定理（\cref{cor:b3:topology:heineborel}(b)）。
\end{exercise}

\begin{exercise}[$\star\star$]\label{exo:b3:topology:8}
(a) 证明 $GL_n(\R)$ 是 $M_n(\R)$ 的开稠密子集，
且它不连通：行列式的符号
将它分成（至少）两块。另一方面，证明
$GL_n(\C)$ 是\omterm{def:b3:topology:pathconnected}{道路}连通的。
\emph{（对 $A, B \in GL_n(\C)$：$\lambda \mapsto \det\bigl((1 -
\lambda)A + \lambda B\bigr)$ 是多项式，不恒为零，
故它在 $\C$ 中只有有限个根：取一条
$\lambda$ 在 $\C$ 中从 $0$ 到 $1$ 的、避开这些根的路径。）}
(b) 证明稠密性：$A + \varepsilon I$ 在
$\varepsilon > 0$ 充分小时可逆。
\end{exercise}

\begin{exercise}[$\star\star$]\label{exo:b3:topology:9}
(a) 证明 $\Q$ 的\omterm{def:b3:topology:components}{连通分支}都是
单点集，并证明 $\Q$ 不是局部紧的（若
$0$ 在 $\Q$ 中的紧\omterm{def:b3:topology:topology}{邻域}将包含 $[-r, r]\cap\Q$，
而其在 $\Q$ 中的\omterm{def:b3:topology:interior}{闭包}不紧：在无理数处截断即可。）
(b) 补全\cref{ex:b3:topology:sinecurve}：不存在\omterm{def:b3:topology:pathconnected}{道路}连接
$(0,0)$ 与 $\Gamma$。\emph{（若 $\gamma = (\gamma_1, \gamma_2)$
是这样的\omterm{def:b3:topology:pathconnected}{道路}，令 $t^* = \sup\{t : \gamma_1(t) = 0\}$；当 $t
> t^*$ 时点在图像上移动；取 $t_k \downarrow
t^*$，使 $\gamma_1(t_k) \to 0$，且所取横坐标使 $\sin$
交替取值 $\pm1$——由介值定理提供这些点——
与 $\gamma_2$ 在
$t^*$ 处的连续性矛盾。）}
\end{exercise}

\begin{exercise}[$\star\star\star$]\label{exo:b3:topology:10}
中三分康托尔集 $C = \bigcap_n C_n$，其中 $C_0 =
[0,1]$，$C_{n+1}$ 是从
$C_n$ 的每个区间中去掉其开的中间三分之一所得。
(a) 证明 $C = \bigl\{\sum_{n\geq1} a_n3^{-n} : a_n \in
\{0,2\}\bigr\}$，且 $C$ 是紧的、\omterm{def:b3:topology:interior}{内部}为空，
且没有孤立点（\emph{完全集}）。
(b) 证明 $C$ 是完全不连通的。
(c) 证明 $x \mapsto (a_n(x))_n$ 是 $C \to
\{0,2\}^\N$ 的\omterm{def:b3:topology:continuity}{同胚}（赋予离散两点
空间的\omterm{def:b3:topology:constructions}{乘积拓扑}）；由此推出 $C$ 不可数，且 $C \times C
\cong C$。
\end{exercise}

\begin{exercise}[$\star\star\star$]\label{exo:b3:topology:11}
球极投影：从北极 $N = (0, \dots,
0, 1)$ 出发，$S^n \subseteq \R^{n+1}$ 上的映射 $\sigma(x) =
\frac{(x_1, \dots, x_n)}{1 - x_{n+1}}$ 是 $S^n
\setminus \{N\} \to \R^n$ 的\omterm{def:b3:topology:continuity}{同胚}（显式给出其逆）。由此推出
$S^n$ \omterm{def:b3:topology:continuity}{同胚}于\omterm{def:b3:topology:locallycompact}{单点紧化}
$\widehat{\R^n}$，且去掉 $S^n$ 中\emph{任意}一点
后所得空间\omterm{def:b3:topology:continuity}{同胚}于 $\R^n$。
\end{exercise}

\begin{exercise}[$\star\star\star$]\label{exo:b3:topology:12}
（\omterm{def:b3:topology:topology}{拓扑}学家正弦曲线）设
\[
T = \bigl\{(x, \sin\tfrac1x) : 0 < x \leq 1\bigr\}
\cup \bigl(\{0\}\times\intcc{-1}1\bigr) \subseteq \R^2 .
\]
(a) 证明 $T$ 是紧的，且它是
图像部分的\omterm{def:b3:topology:interior}{闭包}。
(b) 证明 $T$ 是连通的\emph{（图像作为\omterm{def:b3:topology:continuity}{连续映射}的像是连通的，
其\omterm{def:b3:topology:interior}{闭包}仍连通）}。
(c) 证明 $T$ \emph{不是}\omterm{def:b3:topology:pathconnected}{道路}连通的：不存在连续
\omterm{def:b3:topology:pathconnected}{道路}连接 $(0,0)$ 与 $(1, \sin1)$。\emph{（若 $\gamma =
(\gamma_1, \gamma_2)$ 是这样的\omterm{def:b3:topology:pathconnected}{道路}，令 $t^* = \sup\{t :
\gamma_1(t) = 0\}$；在 $t^*$ 之后，$\gamma_1$ 取遍
充分小的正值（由介值定理），于是
$\gamma_2$ 在 $\pm1$ 之间振荡，对每个区间
$(t^*, t^* + \delta)$ 都成立，与 $t^*$ 处的连续性矛盾。）}
(d) 由此推出\omterm{def:b3:topology:pathconnected}{道路}连通性严格强于
连通性，并说明这类例子不可能是
$\R^n$ 中的\omterm{def:b3:topology:topology}{开集}：$\R^n$ 的一个\emph{开}连通子集
是\omterm{def:b3:topology:pathconnected}{道路}连通的\emph{（与某个基点可用\omterm{def:b3:topology:pathconnected}{道路}相连的点构成的集合
在其定义域中既开又闭）}。
\end{exercise}

\section{问题：康托尔集与空间填充曲线}

\begin{problem}[{周末问题——皮亚诺曲线存在，以及
它们为何不是同胚}]\label{pb:b3:topology:1}
1890 年，皮亚诺以一个\emph{连续
满射} $[0,1] \to [0,1]^2$ 震惊了分析界：一条填满正方形的曲线。
下面仅用基本工具，借助康托尔集 $C$
（见\cref{exo:b3:topology:10}，其结论可任意使用），构造一条这样的曲线，再证明
这样的映射不可能是单射：正方形
不是曲线。全文中，$C$ 的元素写成 $x =
\sum_{n \geq 1} \frac{2b_n(x)}{3^n}$ 的形式，其中数字 $b_n(x) \in
\{0, 1\}$。

\medskip
\noindent\textbf{第一部分——连续地读取数字。}
\begin{enumerate}
  \item 证明：若 $x, y \in C$ 满足 $\abs{x - y} <
        3^{-m}$，则 $b_n(x) = b_n(y)$ 对所有 $n \leq m$ 都成立。
        \emph{（若第一个不同的数字为 $b_k$，则
        $\abs{x - y} \geq \frac{2}{3^k} - \sum_{j>k}\frac2{3^j}
        = \frac1{3^k}$。）}
  \item 由此推出每个数字函数 $b_n \colon C \to \{0,
        1\}$ 连续，并重新推出\omterm{def:b3:topology:continuity}{同胚} $C
        \cong \{0,1\}^\N$，见\cref{exo:b3:topology:10}(c)。
\end{enumerate}

\medskip
\noindent\textbf{第二部分——一个连续满射 $C \to
[0,1]^2$。}
\begin{enumerate}[resume]
  \item 定义 $g \colon C \to [0,1]$，使 $g(x) = \sum_{n\geq1}
        b_n(x)\,2^{-n}$（把康托尔数字读作二进制）。
        证明 $g$ 连续（用第 1 小题），且
        为满射。它是单射吗？
  \item 定义 $\Phi \colon C \to [0,1]^2$ 如下
        \[
        \Phi(x) = \Bigl(\ \sum_{n\geq1} b_{2n-1}(x)\,2^{-n},\
        \sum_{n\geq1} b_{2n}(x)\,2^{-n}\ \Bigr)
        \]
        （奇数位数字给出横坐标，偶数位
        数字给出纵坐标）。证明 $\Phi$ 连续且
        是满射。
\end{enumerate}

\medskip
\noindent\textbf{第三部分——填补空隙：皮亚诺
曲线。}
\begin{enumerate}[resume]
  \item 集合 $[0,1] \setminus C$ 是可数个
        互不相交的开区间 $(u, v)$（即被挖去的中间
        三分之一段），其端点属于 $C$。定义映射 $f \colon
        [0,1] \to [0,1]^2$，其中 $f = \Phi$ 在 $C$ 上成立，并
        在每个缺口上\emph{仿射地}延拓：
        \[
        f(t) = \frac{v - t}{v - u}\,\Phi(u) +
        \frac{t - u}{v - u}\,\Phi(v),
        \qquad t \in (u, v).
        \]
        证明 $f$ 良定义，且满射到
        $[0,1]^2$ 上。
  \item 证明 $f$ 在 $[0,1]
        \setminus C$ 的每一点处连续（局部仿射），并在
        $C$ 的每一点处连续：给定 $\varepsilon$，取 $m$ 使 $2^{-m} <
        \varepsilon/2$，并利用第 1 小题控制 $\Phi$ 在
        $C$ 上 $x$ 附近的取值，再验证仿射
        插值不会逸出：对于缺口 $(u,v) \subseteq
        (x - 3^{-m}, x + 3^{-m})$ 内，$f(t)$ 的值落在
        线段 $[\Phi(u), \Phi(v)]$ 上，两端都靠近
        $\Phi(x)$。结论：$f$ 是\emph{连续
        满射} $[0,1] \to [0,1]^2$。
  \item 推出 $[0,1] \to [0,1]^n$ 的连续满射
        对每个 $n \geq 2$ 成立，并推出 $\R \to \R^2$ 的连续满射。
\end{enumerate}

\medskip
\noindent\textbf{第四部分——但永不单射。}
\begin{enumerate}[resume]
  \item 证明：若连续\emph{单射} $f \colon [0,1]
        \to [0,1]^2$，则它是到其像上的\omterm{def:b3:topology:continuity}{同胚}
        （\cref{cor:b3:topology:compacthomeo}）。
  \item 证明 $[0,1]$ 与 $[0,1]^2$ 不\omterm{def:b3:topology:continuity}{同胚}：
        适当地去掉一点，并比较连通性
        （\cref{met:b3:topology:connectargument}）。
  \item 结论：连续满射 $[0,1] \to [0,1]^2$
        绝不可能是单射——若为单射，则
        $[0,1]^2 = f([0,1])$ 与 $[0,1]$ \omterm{def:b3:topology:continuity}{同胚}（由第 8 小题），
        与第 9 小题矛盾。论证中，
        $[0,1]$ 的紧性究竟在何处起作用？
  \item （总结）归纳其寓意：存在连续
        满射，但不存在连续双射 $[0,1] \to
        [0,1]^2$。这对“维数”作为\omterm{def:b3:topology:topology}{拓扑}概念意味着什么？
        请准确叙述本题证明的一条定理，并叙述一条合理的、
        但仍超出我们工具范围的命题
        （区域不变性）。
\end{enumerate}

\medskip
\noindent\textbf{第五部分——康托集的算术。}
\begin{enumerate}[resume]
  \item 证明 $C + C = [0, 2]$：给定 $t \in [0, 1]$，
        写成 $t = \sum_n d_n3^{-n}$，其中数字 $d_n \in \{0,
        1, 2\}$，并把每个数字拆分为 $d_n = a_n + b_n$，其中
        $a_n, b_n \in \{0, 1\}$；由此得出 $\frac C2 +
        \frac C2 = [0, 1]$ \emph{（就三进制数字而言，$[0,1]$ 的哪个子集是
        $\frac C2$？）}，然后再作伸缩。
        永远不需要进位——请说明这为何是关键所在。
  \item 从几何上解释：“康托尘” $C \times
        C \subseteq \R^2$（\omterm{def:b3:topology:interior}{内部}为空）在对角线上投下
        \emph{完整}的影子：投影 $(x,
        y) \mapsto x + y$ 将其满射到 $[0, 2]$。另外记下
        自相似性 $C = \frac C3 \cup \bigl(\frac C3 +
        \frac23\bigr)$，这正是每幅
        $C$ 示意图背后的方程。
  \item 证明数字交错定义了\omterm{def:b3:topology:continuity}{同胚} $C
        \to C \times C$，并推出 $C \cong C^n$（对每个 $n$）
——康托集是它自己的平方、立方、\dots{}
        哪些熟悉的空间也具有这一性质？
  \item 证明 $\frac14 \in C$ \emph{（计算其三进制
        展开式：$\frac14 = \sum_{k\geq1} 2\cdot3^{-2k}$）}
        尽管 $\frac14$ 不是任何被挖去区间的端点，
        但由计数端点——可推出端点
        构成 $C$ 的一个可数真子集。
  \item 证明第 3 小题中的二进读数映射 $g$ 至多是
        $2$ 对 $1$ 的，并确切描述
        $[0,1]$ 中哪些点恰有两个原像。（$g$ 把 $C$ 压到
        $[0,1]$ 上，通过粘合可数多对点：
        康托阶梯的组合影子，在
        \cref{ch:b3:measure}中再次出现。）
\end{enumerate}

\medskip
\noindent\textbf{第六部分——完美紧集：到处都是康托
集。}非空紧度量空间 $P$ 称为
\emph{完美的}，若它没有孤立点。
\begin{enumerate}[resume]
  \item 设 $P$ 为完美紧集，并设 $x \in P$，$r > 0$。
        证明球 $B(x, r)$ 中包含两点 $y_0
        \neq y_1$ 属于 $P$，从而可在两个\emph{不交}的闭球
        $\bar B(y_0, \rho)$、$\bar B(y_1, \rho)$，它们位于
        $B(x, r)$ 内，以 $P$ 中的点为中心，半径 $\rho
        \leq r/4$，且可任意小。请说明：正是完美性
        （即没有孤立点）恰好使这种
        分割能够在两个新球中\emph{每一个}\omterm{def:b3:topology:interior}{内部}
        重复进行。
  \item 迭代：构造闭集 $F_w \neq \varnothing$
        以有限二进字 $w$ 为指标，使 $F_{w0},
        F_{w1} \subseteq F_w$ 互不相交，且
        $\operatorname{diam}F_w \leq 2^{-\abs w}
        \operatorname{diam}P$。证明对每个无限字
        $\varepsilon \in \{0,1\}^\N$，其交集
        $\bigcap_mF_{\varepsilon_1\cdots\varepsilon_m}$ 为
        单点 $h(\varepsilon)$，\emph{（利用
        紧集 $P$ 的有限交性质）}。
  \item 证明 $h\colon\{0,1\}^\N \to P$ 是单射且
        连续，并由此推出：\emph{每个完美紧度量空间
        都是不可数的}——事实上其基数至少与
        $\{0,1\}^\N$ 的基数一样大。由此得到：$[0,1]$ 与 $C$
        都是不可数的。
  \item 证明 $h$ 是到其像上的\omterm{def:b3:topology:continuity}{同胚}
        \emph{（来自\omterm{def:b3:topology:compact}{紧空间}的连续单射，
        \cref{cor:b3:topology:compacthomeo}）}：每个完美
        紧度量空间都包含一个与
        康托集\omterm{def:b3:topology:continuity}{同胚}的副本。康托集并非怪例，而是
        紧完美性的普适胚芽。
  \item 由此推出：每个可数紧度量空间都有一个
        孤立点，并给出一例，其中孤立点
        稠密但并非全部点：$K = \{0\} \cup
        \{\frac1n : n \geq 1\}$。
  \item （$\R$ 上的康托--本迪克松）设 $F \subseteq \R$ 为
        闭集。称 $x$ 为 $F$ 的\emph{凝聚点}，若
        $x$ 的每个\omterm{def:b3:topology:topology}{邻域}与 $F$ 的交都是不可数的。证明
        对不可数闭集，
        $F$ 的凝聚点构成非空完美闭集 $P$，且 $F
        \setminus P$ 是可数的\emph{（用可数个有理区间覆盖
        非凝聚点，所取的有理区间与
        $F$ 的交集可数）}。由此得出：每个
        $\R$ 的闭子集，要么可数，要么具有
        连续统的基数——对闭集而言，连续统
        假设成立。
\end{enumerate}

\medskip
\noindent\textbf{第七部分——尾声：正则性如何，
离完备有多远。}
\begin{enumerate}[resume]
  \item （曲线的赫尔德正则性）设 $\alpha =
        \frac{\ln 2}{2\ln 3} \approx 0.3155$，并注意恒等式
        $3^\alpha = \sqrt2$。利用第 1 题，证明
        下式成立：
        \[
        \norm{\Phi(x) - \Phi(y)}_\infty \leq 2\,\abs{x -
        y}^\alpha \qquad (x, y \in C),
        \]
        随后沿仿射间隙推广该估计：
        证明第 6 题中的皮亚诺曲线 $f$ 满足
        $\norm{f(t) - f(s)}_\infty \leq 6\abs{t - s}^\alpha$ 在
        整个 $[0,1]$ 上成立 \emph{（同一间隙中的点对用插值处理，
        再对一般点对，经由
        $C \cap [t, s]$ 的端点处理）}。
  \item （指数 $\frac12$ 是一道壁垒）证明：不存在
        满射 $F \colon [0,1] \to [0,1]^2$ 能为
        $\alpha$ 阶赫尔德的，且 $\alpha > \frac12$：将 $[0,1]$
        分成 $N$ 个区间，估计其像的直径，
        并数出网格
        $\bigl(\frac iM, \frac jM\bigr)$，$0 \leq i, j \leq M$ 中
        直径 $< \frac1M$ 的集合所能含有的点数；选取
        $M$ 为 $N^\alpha$ 阶，并令 $N \to \infty$。
        确定我们的曲线（$\alpha \approx 0.3155$）
        相对于该壁垒的位置，并不加证明地记录：希尔伯特曲线
        达到临界指数 $\frac12$。
  \item （导集：度量不完备程度）设 $F$
        为 $\R$ 中的闭集，记 $F'$ 为 $F$ 的极限点集
        （为闭子集），并迭代定义 $F^{(0)} = F$，$F^{(i+1)}
        = (F^{(i)})'$。验证
        \[
        K = \{0\} \cup \Bigl\{\frac1m\Bigr\}_{m \geq 1} \cup
        \Bigl\{\frac1m + \frac1n : m \geq 1,\ n > m^2\Bigr\}
        \]
        是可数紧集，且 $K' = \{0\} \cup
        \{\frac1m\}$, $K'' = \{0\}$, $K''' = \varnothing$
        \emph{（验证第 $m$ 个簇位于
        区间 $(\frac1m, \frac1{m-1})$ 中，从而各簇互不交错）}，
        并验证其孤立点
        在 $K$ 中稠密，与第 21 题的预言相符。描述
        用归纳法构造的过程，对每个 $j$，一个可数紧集
        $K_j$，使 $K_j^{(j)} = \{0\}$ 且 $K_j^{(j+1)} =
        \varnothing$，并与
        第 22 题的完备集对比，对其而言导集运算从不改变集合：
        有限秩恰好与完备性相反。
\end{enumerate}
\end{problem}
